\uuid{BEC3}
\titre{Notion d'équation différentielle}
\auteur{bodin ; croizat ; sacre}
\organisation{exo7}
\licence{CC BY-NC-SA 4.0}
\datecreate{2022-03-08}
\niveau{L1}
\module{Analyse}
\chapitre{Equation différentielle}
\sousChapitre{Résolution d'équation différentielle du premier ordre}
\difficulte{1}
\theme{Notion d'équation différentielle}
\exo7qcmid{hilisit.equations-differentielles.013}

\contenu{
\begin{question}
\qclassification{225.01}
On considère la fonction $f:x\mapsto 2 e^{-x}+3$. Parmi les équations différentielles suivantes, quelles sont celles dont $f$ est solution ?
\begin{answers}
  \good{$y'=-y+3$}
  \good{$y'=y-4 e^{-x}-3$}
  \bad{$y'=2y+3$}
  \good{$y'=-2 e^{-x}$}
\end{answers}
\begin{explanations}
Pour vérifier si une fonction \(f\) est solution d'une équation différentielle du premier ordre, on 	remplace \(y\) par \(f(x)\), \(y'\) par \(f'(x)\) et on regarde si l'égalité est vraie pour tout \(x\) (égalité entre fonctions). Ici \(f'(x)=-2 e^{-x}\). Donc \(f'(x)=-f(x)+3=f(x)-4 e^{-x}-3\) pour tout réel \(x\). Par contre \(2f(x)+3\) n'est pas la même fonction que \( f'(x)\).
\end{explanations}
\end{question}
}
